% Generated draft: requires mathematical and source review.
\documentclass[11pt]{article}
\usepackage[margin=25mm]{geometry}
\usepackage{fontspec}
\setmainfont{Latin Modern Roman}
\usepackage{amsmath,amssymb,mathtools,unicode-math}
\setmathfont{Latin Modern Math}
\usepackage{enumitem,needspace,xurl,hyperref}
\setlength{\parindent}{0pt}
\setlength{\parskip}{5pt}
\setlength{\emergencystretch}{3em}
\newcommand{\N}{\mathbb{N}}
\newcommand{\Z}{\mathbb{Z}}
\newcommand{\Q}{\mathbb{Q}}
\newcommand{\R}{\mathbb{R}}
\newcommand{\C}{\mathbb{C}}
\newcommand{\F}{\mathbb{F}}
\newcommand{\A}{\mathbb{A}}
\newcommand{\PP}{\mathbb{P}}
\newcommand{\E}{\mathbb{E}}
\newcommand{\eps}{\varepsilon}
\newcommand{\dd}{\,\mathrm{d}}
\begin{document}
\section*{UnsolvedMath ID: 30002204}
\textbf{Source problem code:} catalogue entry
\subsection{Definitions and mathematical statement}
The source asks for a description of the structure and minimal free resolutions of Orlik--Terao algebras. It does not specify a ground field or all arrangement hypotheses. The following standard central/projective convention makes the notation explicit without resolving those omissions.

Let \(k\) be a field, let \(V\) be a finite-dimensional \(k\)-vector space, and let \(A=\{H_1,\ldots,H_n\}\) be a finite central hyperplane arrangement with defining nonzero linear forms \(\alpha_i\in V^*\), so that \(H_i=\ker\alpha_i\). Set \(R=k[y_1,\ldots,y_n]\), with \(\deg y_i=1\). For each nonzero linear relation \(c=(c_1,\ldots,c_n)\) satisfying \(\sum_{i=1}^n c_i\alpha_i=0\), write \(S(c)=\{i:c_i\ne0\}\) and define
\[
\partial(c)=\sum_{i\in S(c)}c_i\prod_{j\in S(c)\setminus\{i\}}y_j.
\]
The Orlik--Terao ideal and algebra are
\[
I_A=(\partial(c):c\ne0,\ \textstyle\sum_i c_i\alpha_i=0)\subseteq R,
\qquad
\operatorname{OT}(A)=R/I_A.
\]
For a projective arrangement, this convention uses the corresponding central arrangement of defining linear forms. In particular, the source's notation \(A\subseteq\mathbb P^2\) is understood as an arrangement of projective lines, represented by linear forms on \(k^3\), rather than as a set of points.

A minimal graded free resolution of \(\operatorname{OT}(A)\) over \(R\) is an exact sequence
\[
\cdots\longrightarrow F_2\xrightarrow{d_2}F_1\xrightarrow{d_1}F_0
\longrightarrow R/I_A\longrightarrow0,
\qquad
F_p=\bigoplus_{q\in\Z}R(-q)^{\beta_{p,q}},
\]
whose differentials are homogeneous of degree zero and satisfy \(d_p(F_p)\subseteq(y_1,\ldots,y_n)F_{p-1}\) for \(p\ge1\). Here \(R(-q)\) denotes the grading shift with its generator in degree \(q\), and \(\beta_{p,q}\) are the graded Betti numbers.

\textbf{Problem.} Describe the phenomena exhibited by \(\operatorname{OT}(A)\) in higher dimensions and for reflection arrangements, and determine its minimal graded free resolution over \(R\), including in the case of projective line arrangements in \(\mathbb P^2\). A reflection arrangement is, in the usual finite linear-group setting, the collection of hyperplanes fixed pointwise by reflections, where a reflection is a nonidentity linear transformation whose fixed subspace has codimension one.

The source does not specify the field, characteristic, class of reflection groups, meaning of ``higher dimensions,'' or particular structural invariants meant by ``phenomena.'' Consequently, the structural part remains a research question rather than a uniquely specified proposition, and the displayed conventions do not assert that every field or arrangement class is intended.
\subsection{Short English statement}
Describe the structure of Orlik--Terao algebras in higher dimensions and for reflection arrangements, and determine their minimal graded free resolutions, even for arrangements of lines in \(\mathbb P^2\). The source leaves the ground field, reflection-group class, and precise structural targets unspecified. Its supplied status is partially solved; this draft does not independently verify that assessment.
\subsection{Sources}
\begin{itemize}
\item Local UnsolvedMath registry: supplied record with permanent ID \(30002204\), title \emph{Resolutions and Higher-Dimensional Structure of Orlik--Terao Algebras}, source problem code \texttt{catalogue entry}, in the Oberwolfach Reports Open Problems collection. No individual proposer is supplied.
\item Supplied report citation: \emph{Mini-Workshop: Cohomology Rings and Fundamental Groups of Hyperplane Arrangements, Wonderful Compactifications, and Real Toric Varieties} (2013), Oberwolfach Reports, pp.\ 2939--2983, DOI \texttt{10.4171/owr/2012/49}. Original OWR problem: \url{https://doi.org/10.4171/owr/2012/49}. The record separately gives proposed year 2012; the relationship between that year and the citation's 2013 label is not resolved here.
\item Supplied literature source, labeled ``Orlik--Terao blowup algebra'': \url{https://arxiv.org/abs/1701.03470}.
\item Supplied literature source, labeled ``Orlik--Terao ideal and relation space'': \url{https://doi.org/10.1016/j.aam.2015.08.003}.
\item Literature assessment dated exactly \texttt{2026-08-21}: the supplied curation labels the problem partially solved and reports known Cohen--Macaulayness, blowup descriptions, modular decompositions, and important special classes, while reporting no general minimal free resolution or complete higher-dimensional and reflection-arrangement theory. This is supplied dated curation, not an independently verified literature assessment; its status and cited claims require independent verification.
\item The registry labels the statement \texttt{corrected\_verified} and reports that its mathematical content was checked against the official Oberwolfach report, with notation, context, and question form normalized without changing the source-level claim. This is supplied verification metadata, not a verification performed for this draft.
\end{itemize}
\end{document}
